Long-context Transformer inference stores a key and value for every visible token, so memory grows linearly with context length. Existing responses usually address only one part of the problem. Quantization reduces storage but not the width of attention. Eviction and retrieval reduce attention work, but typically discard history without a formal local error guarantee. Sequential residual coding can compress historical states effectively while making random access expensive because one block depends on the reconstruction of many earlier blocks. RACK-KV is designed around this three-way tension. It keeps a short recent window exact, compresses older keys and values into independently decodable block-local anchor/residual codes, stores compact geometric summaries for each block, and skips a historical contribution only when a query-dependent rigorous certificate proves that the resulting local reconstructed attention error is below a chosen tolerance. The certificate is evaluated with directed-rounding multiprecision arithmetic, so floating-point uncertainty cannot authorize an unsafe skip. In a representative-layer study derived from a fixed \modeldisplay{} checkpoint, RACK-KV reduced serialized KV storage by about \num{36.5}\% while maintaining mean attention-output error \num{2.29e-3}. In a verified full 32-layer teacher-forced evaluation on two 128-token prompt categories, mean negative log-likelihood remained within \num{0.002} of the Full-KV reference and top-1 next-token agreement stayed above \num{98}\% for both RACK-KV streams. No certificate violation was observed, but certified skip utilization remained low and no physical block decodes were avoided. RACK-KV therefore provides evidence that compression, independent random access, and rigorous local skipping can coexist in one design, while broader prompt coverage, stronger end-to-end theory, and real systems acceleration remain future work.
